\section{Benchmarks of the reference implementation}\label{app:benchmarks}

The main text quotes the results of several numerical experiments in the
reference implementation; their details are collected here. Every number is
produced by code that accompanies the site and is checked by its tests.

\paragraph{Separation with noise.} The wave and rotor are those of
Section~\ref{sec:one-pixel}. With $\tE$ and $\tB$ at one pixel the near field
pays the same $\sqrt2$ as the wave in the two components it shares with the
wave, and nothing in the other three. An array that reads only $\tE$ fits the
rotor as a point quadrupole of unknown moment and position; for a pixel and six
neighbours at distance $b$ along three axes and $b \ll r$, the factor on the
wave's error lies between $0.22\,r/b$ and $0.55\,r/b$ over the rotor's
orientations, and the array overtakes the magnetic part beyond $b = 0.29\,r$.
In return the wave costs the rotor's moment and position factors of $1.5$ and
$1.3$ from $\tE$ alone, and $1.07$ and $1.04$ with $\tB$ too. Modelling the
rotor as a point leaves a false wave of $0.005$ of the true one, from its next
multipole, which matters only above $\rho \approx 270$; ignored altogether, the
rotor is read as $70$ times the wave by the array.

\paragraph{Depth with noise.} For a person $3\,\mathrm{m}$ away and neighbours
$0.5\,\mathrm{m}$ out, the bound of Eq.~\eqref{eq:depth-bound} on $\sigma_r/r$
is between $0.90/\rho$ and $1.08/\rho$, depending on direction, against
$1.16/\rho$ from the closed form (Fig.~\ref{fig:separation-noise}b): wider
stencils gain from the curvature of the field across them. For a person walking
at $1.4\,\mathrm{m\,s^{-1}}$ the magnetic bound is $\sigma_r/r =
1.2\times10^{8}/\rho$. Bodies on one line: a rod of length $L$ pointing at the
camera appears as a point $L^2/3r$ nearer than its centre; the spread $s$ of the
depths enters the second derivatives as a share $s^2/r^2$, which the seven
pixels read to $1/\rho_H$, with $\rho_H$ between $8.1$ and $10.5$ times
$\rho\,(b/r)^2$ depending on direction, so that two equal masses, $s =
\Delta/2$, are told from one when $\Delta > 2r/\rho_H^{1/2}$, between $0.62$
and $0.70$ times $(r^2/b)/\rho^{1/2}$.

\paragraph{The worked inversion.} The array of Section~\ref{sec:near-field}'s
worked case is a square a metre across, of $11\times11$ pixels $10\,\mathrm{cm}$
apart, centred on the camera and facing the safes; the noise on each component
is set by $\SNR = 2Gm/z^2\sigma a$ for $m = 324\,\mathrm{kg}$, $z =
3.5\,\mathrm{m}$ and spacing $a$. The image is the Gaussian posterior on point
masses at the nodes of a grid $25\,\mathrm{cm}$ apart, $5\,\mathrm{m}\times
2.5\,\mathrm{m}$, in the plane of the safes, each with prior standard deviation
equal to the gold's mass; it is computed through the singular values of the
design matrix, which stays exact at any noise. A safe's mass is the sum over the
cells of its footprint, $0.9\,\mathrm{m}$ wide. Gold $0.2\,\mathrm{m}$ farther
back and $14\%$ heavier, so as to pull as hard on the array's centre, is told
apart at $\SNR \approx 10^3$, from the way the field changes across the array.
The posterior uncertainty is $\pm 730\,\mathrm{kg}$ at $\SNR = 10^3$ and $\pm
490\,\mathrm{kg}$ at $10^6$ for the metre-wide array, and $\pm 350\,\mathrm{kg}$
at $10^6$ for one six metres across; these belong to the array and the noise,
not to the data. Tilting changes nothing in principle, since each pixel reads
the whole tensor wherever it points: turned through $70^\circ$, the metre-wide
array tells the two safes apart $9\%$ better, and the six-metre one, whose near
edge then comes within $0.7\,\mathrm{m}$ of the safes, thirty times better. The
positive reconstruction is non-negative least squares on the same grid, with no
prior, over forty draws of the noise.

\paragraph{The trembling lattice.} The ball and the shell of
Fig.~\ref{fig:tremor} are the sites of a simple cubic lattice of unit spacing
within a radius of six spacings, the shell keeping only those beyond four, each
site joined to its six nearest neighbours by springs of unit stiffness and to
its twelve next-nearest by springs half as stiff, which makes a stable elastic
solid. The masses are equal within each body and add to the same total in
both. The normal modes give the response of every moment directly: a mode of
frequency $\omega_n$ moves $q_h$ by $c_n$ per unit of its mass-weighted
amplitude, and in equilibrium $\langle\dot q_h^2\rangle = k_BT\sum_n c_n^2$ and
$\langle q_h^2\rangle = k_BT\sum_n c_n^2/\omega_n^2$. The first sum equals
$k_BT\int\rho\,|\nabla h|^2\,\dd V$ to rounding error for any stiffnesses, as
Eq.~\eqref{eq:thermal-moment} requires; for $h = x$ all of it lies in the free
translation, the Brownian motion of Eq.~\eqref{eq:thermal-multipole} with
$\ell = 1$. The steel ball's numbers use Rayleigh's method with the
displacement $\mathbf{u} = (q/R)(-x/2, -y/2, z)$, which has no divergence and
moves only the surface: effective mass $3M/10$, stiffness $3\mu V/R^2$, a
quadrupole $\tfrac15 MRq\,\mathrm{diag}(-1,-1,2)$ and so a tide
$36GMRq/5r^5$ on its axis. Rayleigh's method overestimates the frequency;
these are estimates, not the exact modes of a steel sphere.

For a quantum system the jitter is no longer $k_BT\chi(0)$. With the
static compliance $\chi_0$ and the sum $A$ of Eq.~\eqref{eq:sum-rule} both
finite, the measure $(2/\pi\chi_0)\,\mathrm{Im}\,\chi(\omega)\,\dd\omega/\omega$
has unit weight and mean square frequency $A/\chi_0$, and the quantum
fluctuation--dissipation theorem gives the jitter as its average of
$\tfrac12\hbar\chi_0\,\omega\coth(\hbar\omega/2k_BT)$. Since
$\sqrt{x}\coth(b\sqrt{x})$ is concave in $x$,
\begin{equation}\label{eq:jitter-bounds}
k_BT\chi_0 \le \langle\delta q^2\rangle \le
\frac{\hbar}{2}\sqrt{A\chi_0}\,\coth\!\Big(\frac{\hbar}{2k_BT}\sqrt{\frac{A}{\chi_0}}\Big),
\end{equation}
with the upper bound reached by a single mode. The reference implementation
checks both sides on random spectra, and checks that the shell drive of
Eq.~\eqref{eq:thermometer} has no monopole or dipole, interacts with point
masses placed at random outside it exactly as $-f\,D\!:\!\tE$, and makes the
field $6GfD/a^5$ at its centre.
